\chapter{Curvature ratios and local translator rigidity}\label{sec:translator-curvature-rigidity}
We use the translator convention of Definition \ref{def:translator-jacobi-weight}. The primary source is \cite{MSHS14}, Lemma 2.1, pp.~5--7, and the calculation in Theorem A, pp.~10--12. We state local consequences of a local maximum separately from a complete global classification. No analyticity or unique-continuation theorem is implicit in that separation.

\begin{proposition}[Height Hessian and the translator Simons equation]\label{prop:translator-tensor-simons}
On a smooth two-sided translator, with $h=\langle X,v\rangle$,
\[
 \nabla^2h=HA,\qquad \nabla H=-A(v^T),\qquad
 \Delta A+\nabla_{v^T}A=-|A|^2A.
\]
Here the Laplacian on $A$ is the rough covariant Laplacian. Consequently
\[
 \mathcal D_v|A|^2=2|\nabla A|^2-2|A|^4,
 \qquad \mathcal D_vH^2=2|\nabla H|^2-2|A|^2H^2.
\]
\end{proposition}
\begin{proof}
The ambient Hessian of the linear function $\langle x,v\rangle$ vanishes. Its intrinsic Hessian is the pairing of $v$ with the vector second fundamental form $-A\nu$, which is $HA$ because $\langle v,\nu\rangle=-H$. Differentiating this latter scalar equation gives $\nabla H=-A(v^T)$.

At a point in a geodesic frame, differentiate the gradient equation and use Codazzi to move the derivative of $A$:
\[
 (\nabla^2H)_{ij}=-(\nabla_{v^T}A)_{ij}-H(A^2)_{ij}.
\]
The second term uses $\nabla v^T=\nabla^2h=HA$. Substitute into the Euclidean case of Theorem \ref{thm:spaceform-simons}; the two $HA^2$ terms cancel, yielding the tensor equation. Metric compatibility and the product rule give its norm equation. Either tracing the tensor equation or using Proposition \ref{prop:translator-angle-jacobi} gives $\mathcal D_vH=-|A|^2H$, and the product rule gives the final formula. These are pointwise identities without completeness or sign assumptions on $H$.
\end{proof}

\begin{proposition}[The normalized shape tensor and curvature ratio]\label{prop:translator-curvature-ratio}
On an open region where $H\ne0$, let $B=A/H$ and $R=|A|^2/H^2=|B|^2$. Then
\[
 \Delta B+\nabla_{v^T}B+2\nabla_{\nabla\log|H|}B=0,
\]
and
\[
 \Delta R+\langle v^T+2\nabla\log|H|,\nabla R\rangle
 =2|\nabla B|^2
 =\frac2{H^4}|H\nabla A-\nabla H\otimes A|^2.
\]
In particular the ratio is subharmonic for the displayed drift operator on that region.
\end{proposition}
\begin{proof}
Insert $A=HB$ into Proposition \ref{prop:translator-tensor-simons}. The product rule gives
\[
 -|A|^2HB=H(\Delta B+\nabla_{v^T}B)
 +(\mathcal D_vH)B+2\nabla_{\nabla H}B.
\]
The terms with $\mathcal D_vH=-|A|^2H$ cancel. Divide by $H$, which is legitimate on the stated region, to get the tensor equation. Applying its scalar drift operator to the squared norm gives twice the squared covariant derivative. Finally $\nabla B=(H\nabla A-\nabla H\otimes A)/H^2$. This proves the coefficient and the nonnegative right side in the curvature-ratio equation.
\end{proof}

\begin{proposition}[What a maximum of the ratio implies]\label{prop:translator-ratio-maximum}
If $R$ has a local maximum at a point where $H\ne0$, there is a connected neighborhood on which $R$ is constant and $B=A/H$ is parallel. If $R$ attains its global maximum on a connected component of $\{H\ne0\}$, the same conclusion holds throughout that component.
\end{proposition}
\begin{proof}
Choose a relatively compact connected coordinate neighborhood where $H$ stays away from zero and the local maximum is a maximum on the whole neighborhood. The operator in Proposition \ref{prop:translator-curvature-ratio} is uniformly elliptic there and has bounded smooth drift and zero zeroth-order term. The scalar strong maximum principle gives constancy of $R$ on this neighborhood. Its equation then gives $\nabla B=0$.

For a global maximum on a connected component, the set where that maximum is attained is relatively closed by continuity and is relatively open by the same argument. It is nonempty, hence is the whole component. No uniform ellipticity constants at infinity are needed, since the argument is local at each maximum point. The scalar strong maximum principle is a background analytic input. A local maximum alone is not silently promoted to a global conclusion by that principle.
\end{proof}

\begin{lemma}[A parallel normalized shape tensor has rank one]\label{lem:translator-parallel-shape-rank-one}
Let $U$ be a connected open region of a translator where $H\ne0$ and $B=A/H$ is parallel. Then $B$ is the orthogonal projection onto a one-dimensional parallel distribution. Thus $A$ has rank one on $U$, $|A|^2=H^2$, and the induced metric is flat.
\end{lemma}
\begin{proof}
The symmetric parallel tensor $B$ has constant eigenvalues on connected $U$ and trace one. Codazzi applied to $A=HB$ gives
\[
 (\nabla_iH)B_{jk}=(\nabla_jH)B_{ik}.
\]
If $\nabla H$ is nonzero at a point, choose the first frame vector in its direction. Taking $i=1$ and $j\ge2$ shows $B_{jk}=0$ for all such $j,k$. Symmetry and trace one make $B$ a rank-one orthogonal projection there, and parallel transport gives the same eigenvalues throughout $U$.

If no such point exists, $H$ is a nonzero constant on $U$. But $J_vH=0$ would then give $|A|^2H=0$, contradicting $H=\operatorname{tr}A\ne0$. Thus the preceding case must occur. The nullity and range distributions of a parallel projection are parallel. In the Gauss equation every sectional curvature is a difference of two products of entries of a rank-one form and is zero, so the induced metric is flat.
\end{proof}

\begin{proposition}[A local grim-reaper product]\label{prop:translator-local-grim-product}
Under the hypotheses of Lemma \ref{lem:translator-parallel-shape-rank-one}, every point has a neighborhood whose immersed image, after a rigid motion, is an open piece of a grim-reaper curve times $\mathbb R^{n-1}$. The velocity may have a tangential component in the Euclidean factor. Its nonzero projection onto the grim-reaper plane determines the curve's speed.
\end{proposition}
\begin{proof}
Work in a simply connected small flat neighborhood. Parallel transport gives a parallel orthonormal frame with $e_1$ in the range of $B$ and $e_2,\ldots,e_n$ in its kernel. The existence of such a frame for a flat connection is the background differential-geometric input. Torsion-freeness makes these fields commute, so there are local coordinates $(s,y_2,\ldots,y_n)$ with these coordinate fields. For $\alpha\ge2$, the ambient derivative of $dX(e_\alpha)$ in every tangent direction is zero: its tangential part vanishes by parallelness and its normal part is $-A(\cdot,e_\alpha)\nu=0$. Thus these vectors are constant ambient vectors $a_\alpha$.

Integration gives $X(s,y)=\gamma(s)+\sum_{\alpha\ge2}y_\alpha a_\alpha$. The unit-speed curve $\gamma$ lies in an affine two-plane orthogonal to their span. Projecting the translator equation to that plane shows that it translates with the projected velocity $v_0$. This projection is nonzero, since otherwise its curvature would vanish, contrary to $H\ne0$. Rotate within the plane to write $v_0=ce_2$, $c>0$. The equation $H=-\langle v_0,\nu\rangle\ne0$ makes its vertical normal component nonzero, so the curve is locally a graph over the horizontal axis. Proposition \ref{prop:grim-reaper-speed-family} identifies it as a grim-reaper graph piece. No completeness or global continuation was used.
\end{proof}

\begin{theorem}[Complete rigidity with a global curvature-ratio maximum]\label{thm:translator-complete-ratio-rigidity}
Let $X:\Sigma^n\to\mathbb R^{n+1}$ be connected, complete, boundaryless, smooth and two-sided, with $H$ nowhere zero. Suppose it translates with a fixed velocity $v$ and $|A|^2/H^2$ attains a global maximum. Then its image is a complete grim-reaper curve times an $(n-1)$-dimensional Euclidean space, possibly tilted relative to $v$. This is a restricted complete formulation of the rigidity mechanism in \cite{MSHS14}, Theorem A, pp.~10--12; the explicit hypotheses avoid an unstated continuation across $\{H=0\}$.
\end{theorem}
\begin{proof}
Proposition \ref{prop:translator-ratio-maximum} and Lemma \ref{lem:translator-parallel-shape-rank-one} give a parallel rank-one projection $B$ on all of $\Sigma$. The images under $dX$ of its kernel define one fixed ambient $(n-1)$-plane $L$. Indeed, parallel transport of a kernel vector along any curve has zero ambient derivative, since both its covariant derivative and its second fundamental form with the curve tangent vanish. Connectedness makes this ambient plane independent of the base point.

For any $a\in L$, the unique tangent field $V_a$ with $dX(V_a)=a$ has constant length. It is complete: an integral curve on a finite time interval stays in a bounded intrinsic ball, compact by completeness, so the smooth ODE can be continued. Along its flow $X$ changes by $ta$. Therefore the image is invariant under all translations in $L$.

Project onto $L^\perp$, an affine two-dimensional factor, and let $v_0$ be the projection of $v$. It is nonzero because $H\ne0$. Choose a unit vector $q\in L^\perp$ perpendicular to $v_0$. The projection $q^T$ on the surface never vanishes: otherwise $q$ would be normal, giving $\langle v_0,\nu\rangle=0$ and $H=0$. Thus $E=q^T/|q^T|$ is a global smooth unit tangent field spanning the range of $B$, and its flow exists for all real time by the same bounded-speed argument. Its image, after projection onto $L^\perp$, is a complete unit-speed planar translator with nonzero curvature and everywhere positive $q$ component of its tangent. The graph ODE in Proposition \ref{prop:grim-reaper-speed-family} identifies it with the full grim reaper: any omitted endpoint strictly inside its maximal graph interval would have bounded slope and finite arclength, contradicting continuation of this complete unit-speed flow; the two genuine asymptotes have infinite arclength.

The flows of $E$ and the fields $V_a$ span the tangent space everywhere. Their accessible orbit from any point is therefore open, and the orbits partition the connected manifold into open sets, so there is only one. The $L$ flows do not change the two-plane projection; the $E$ flow stays on the same grim-reaper curve by its ODE uniqueness. Hence no additional image component occurs, and translation invariance in $L$ gives the asserted full product. The conclusion describes the image; it does not rely on a separate analytic unique-continuation theorem.
\end{proof}
